%% LyX 2.5.1 created this file.  For more info, see https://www.lyx.org/.
%% Do not edit unless you really know what you are doing.
\documentclass[12pt,english,preprint,number]{elsarticle}
\usepackage[utf8]{inputenc}
\synctex=-1
\usepackage{color}
\usepackage{colortbl}
\usepackage{babel}
\usepackage{float}
\usepackage{booktabs}
\usepackage{amsmath}
\usepackage{amssymb}
\usepackage[bookmarks=false,
 breaklinks=false,pdfborder={0 0 1},backref=false,colorlinks=true]
 {hyperref}
\hypersetup{
 allcolors=blue}

\makeatletter

%%%%%%%%%%%%%%%%%%%%%%%%%%%%%% LyX specific LaTeX commands.
\newcommand{\lyxmathsym}[1]{\ifmmode\begingroup\def\b@ld{bold}
  \text{\ifx\math@version\b@ld\bfseries\fi#1}\endgroup\else#1\fi}

%% Because html converters don't know tabularnewline
\providecommand{\tabularnewline}{\\}

%%%%%%%%%%%%%%%%%%%%%%%%%%%%%% User specified LaTeX commands.
\PassOptionsToPackage{english}{babel}
\usepackage{babel}


% Useful packages
\usepackage{array}
\usepackage{url}
\usepackage{graphicx}
\graphicspath{{../figures/}{../../wpm/}{../../notebooks/cbed/results/}{../../notebooks/verification/figures/}}
\DeclareGraphicsExtensions{.pdf,.png,.jpg}
\newcommand{\safeincludegraphics}[2][]{%
    \IfFileExists{../figures/#2.pdf}{\includegraphics[#1]{../figures/#2}}{%
        \IfFileExists{../figures/#2.png}{\includegraphics[#1]{../figures/#2}}{%
            \IfFileExists{../figures/#2.jpg}{\includegraphics[#1]{../figures/#2}}{%
                \IfFileExists{../../notebooks/cbed/results/#2.pdf}{\includegraphics[#1]{../../notebooks/cbed/results/#2}}{%
                    \IfFileExists{../../notebooks/cbed/results/#2.png}{\includegraphics[#1]{../../notebooks/cbed/results/#2}}{%
                        \IfFileExists{../../notebooks/verification/figures/#2.pdf}{\includegraphics[#1]{../../notebooks/verification/figures/#2}}{%
                            \IfFileExists{../../notebooks/verification/figures/#2.png}{\includegraphics[#1]{../../notebooks/verification/figures/#2}}{%
                                \fbox{\parbox[c][0.25\textheight][c]{0.9\linewidth}{\centering Missing figure:\\ \texttt{\detokenize{#2}}}}%
                            }%
                        }%
                    }%
                }%
            }%
        }%
    }%
}
\usepackage{tikz}
\usetikzlibrary{positioning}
\newcolumntype{L}[1]{>{\raggedright\arraybackslash}p{#1}}

\journal{Ultramicroscopy}
\biboptions{numbers,sort&compress}



\makeatother

\begin{document}
\begin{frontmatter}
\title{Wide-Angle Multislice Propagation for Electron Microscopy: Comparison
of Fresnel, Angular-Spectrum, and Wave-Propagation Methods}
\author[inst1]{David Landers}
\ead{david.landers@cea.fr}
\author[inst1]{Jean-Luc Rouviere}
\ead{jean-luc.rouviere@cea.fr}

\end{frontmatter}
\
\begin{abstract}
Conventional Fresnel multislice (F-MS) is the standard method for
modeling electron--sample interactions in transmission electron microscopy
(TEM). It is computationally efficient and can simulate both perfect
and defective crystals. However, F-MS loses accuracy at large scattering
angles. Here, we investigate FFT-based alternatives that retain these
practical advantages while improving wide-angle accuracy without the
high computational cost of directly solving the wave equation. We
introduce wave-propagation multislice (WP-MS), a minor adaptation
of the wave propagation method (WPM) from scalar optics, and compare
it with F-MS and angular-spectrum multislice (AS-MS). The three methods
are first benchmarked against a direct ordinary-differential-equation
(ODE) solution of the second-order Klein--Gordon Helmholtz equation
using the on-axis and high-angle beam amplitudes of Au~{[}100{]}.
We then compare their predictions for convergent-beam electron diffraction
(CBED) from Si~{[}111{]} and Au~{[}100{]}. AS-MS removes most of
the wide-angle error of F-MS at essentially the same per-slice computational
cost. WP-MS is closest to the ODE reference, but its additional accuracy
over AS-MS is modest under the conditions studied. We therefore identify
the conditions under which F-MS remains adequate, recommend AS-MS
as the practical wide-angle default, and reserve WP-MS for applications
requiring the most complete one-way split-step model. 
\end{abstract}
\begin{keyword}
electron microscopy \sep multislice \sep electron diffraction \sep
TEM image simulation \sep Fresnel propagator \sep angular spectrum
\sep wave propagation method 
\end{keyword}

\section{Introduction}

\label{sec:introduction}

Electron microscopy (EM) has recently entered a new era. For a long
time, its central role was to produce images of matter, often interpretable
only through comparison with image simulations. In this traditional
workflow, simulations were used as a forward model: a structural hypothesis
was proposed, an image was calculated from this structure, and the
agreement with experiment was improved by trial and error. In this
sense, image simulation has always been an indispensable tool for
extracting reliable structural information from EM data.

This situation is now changing. In modern approaches, especially ptychography
and related phase-retrieval methods, the recorded data are no longer
simply interpreted as images. Instead, they are used as input for
numerical reconstruction algorithms that aim to recover the sample
structure, projected potential, or phase shift. In these inverse methods,
the electron--matter interaction model is not merely used after the
experiment for comparison; it becomes part of the reconstruction itself.
Consequently, the accuracy of the forward model is more critical than
ever.

A complete description of electron propagation through matter would
require solving a many-body quantum-mechanical problem including elastic
and inelastic scattering, multiple scattering, and, in principle,
backscattering processes. Such a brute-force treatment is not feasible
for realistic specimens of interest in transmission electron microscopy
(TEM). Practical and accurate propagation models are therefore essential.
In this work, we focus on multislice approaches, leaving aside Bloch-wave
methods, which remain powerful for periodic crystals ~\citet{Bethe_theorie_1928,Hirsch1965}
but are less general for defects, interfaces, amorphous regions, finite
objects, and iterative reconstruction schemes.

The multislice algorithm introduced by Cowley and Moodie ~\citet{cowley1957scattering,ishizuka1977practical}
provides such a practical framework. It was originally motivated by
a physical-optics picture in which the specimen is represented as
a succession of thin phase objects, separated by free-space propagation.
In modern terms, the electron wave is alternately multiplied by a
transmission function associated with the projected potential of a
slice and propagated to the next slice. This simple and efficient
construction has become one of the standard tools for simulating electron
scattering in TEM.

The conventional implementation of multislice (F-MS) uses a Fresnel,
or paraxial, propagator between slices. - - - - - - - - I insert back
David's review here - - - -

F-MS is usually adequate for conventional high-energy TEM and STEM
simulations of thin, low-$Z$ specimens, where other uncertainties
such as thermal diffuse scattering, detector response, or microscope
aberration characterization often dominate the comparison with experiment.
However, for low incident electron energies, thick specimens, or high-$Z$
materials, F-MS can exhibit significant errors because of its paraxial
approximation and omission of backscattering ~\citet{RotherScheerschmidt2009,hillier_role_2017,WackerSchroeder2015,ming2013pcms}.
AS-MS is the minimal FFT-based correction to the paraxial propagation
step, but it still treats the specimen as a sequence of thin transmission
functions separated by vacuum propagation. More complete wide-angle
methods therefore seek to improve the propagation inside the specimen,
or to use the full second-order wave equation as a reference.

Further improvements in accuracy, particularly the treatment of large-angle
scattering and backward-propagating components, incur significant
computational cost. A more accurate electron microscopy multislice
algorithm is the Chen--Van Dyck method, also known as Fully Corrected
Multislice (FC-MS) \citet{ming2013pcms,ChenVanDyck1997}. The complexity
of FC-MS arises because the full second-order wave equation contains
a square-root operator that mixes a transverse derivative term with
the potential and cannot be evaluated as a single Fourier-space propagator.

FC-MS addresses this problem by expanding the square-root operator
in a Taylor series and retaining higher-order correction terms that
couple the potential to transverse derivatives. In this way the method
accounts more accurately for how the wave evolves within the potential
of each slice, rather than treating the slice only as a projected
phase grating followed by free-space propagation. This allows improved
treatment of larger scattering angles, low-energy electrons, thicker
specimens, and high-$Z$ materials and, in its more complete formulation,
also gives access to backward-propagating components.

5 years before the FC-MS was introduced in the field of scalar light
optics, a solution to the identical mathematical problem of solving
a Time Independent Wave Equation to higher accuracy than the F-MS
or AS-MS method had already been proposed, known as the Padé Beam-Propagation
Method (Padé BPM) ~\citet{Hadley1992}. In light optics, the time-independent
Helmholtz equation also leads to a square-root propagation operator
when it is factorized into forward- and backward-propagating branches.
Padé beam-propagation methods approximate this square-root operator
using rational functions rather than the Taylor polynomials proposed
by Chen and Van Dyck. Because a Padé approximant is a ratio of polynomials,
it can represent the square-root dispersion relation more accurately
over a wider angular range than a truncated Taylor series of the same
order. In practice, this often gives improved stability and accuracy
for large-angle propagation, allowing lower-order Padé approximants
to achieve accuracy that would otherwise require higher-order Taylor
corrections. In either case, approximating the square-root operator
remains computationally expensive and incurs significant overhead
relative to F-MS or AS-MS.

In electron optics, a still more accurate reference can be obtained
by directly integrating the full second-order time-independent wave
equation as an ordinary differential equation (ODE). Numerous authors
have used such solutions as references for assessing the accuracy
of F-MS and the FC-MS method. Wacker and Schröder provide a clear
derivation of the ODE and operator-approximation methods and benchmark
them against conventional F-MS and direct Runge--Kutta integration
of the second-order Schrödinger equation. Their results again showed
that the F-MS method can introduce substantial errors at low voltages,
large specimen thicknesses, and high scattering angles. Rother and
Scheerschmidt approached the problem from a relativistic perspective,
using ODE solutions of the Dirac, Klein--Gordon, and relativistically
corrected Schrödinger wave equations and comparing them with a paraxial
multislice solution. Their work similarly indicates that the dominant
high-angle error arises from the parabolic propagation step used in
F-MS. Hillier et al.~\citet{hillier_role_2017} show also that the
second order derivative term is important to track the positions of
Laue Zone rings in low energy CBED patterns of Silicon at 20keV. However,
out of all the methods discussed so far, it is only \citet{ming2013pcms}
that attempted to account for high angle effects in a way that can
exploit the Fourier transform, with their spherical propagation method.
The other series expansion and ODE methods cannot exploit the FFT
as directly as split-step multislice, limiting their practical use
as a forward model within inverse reconstruction algorithms.

The light-optics community has sought alternatives that retain the
computational efficiency of the FFT while avoiding the numerical difficulties
of evaluating Padé BPM approximants and still providing a nonparaxial
solution to the time-independent wave equation. One such solution
to this problem is the Wave Propagation Method (WPM), introduced by
Brenner in 1993~\citet{Brenner1993}. In the WPM, the potential or
refractive index medium is treated as locally homogeneous over small
transverse regions, and a corresponding homogeneous one-way propagator
is applied for each representative refractive index (see \ref{fig:wpm_explainer}
for a visual explanation of the method). The resulting fields are
then recombined according to the local refractive-index distribution.
This retains the potential-dependent wavelength within the propagation
step without explicitly forming the full dense square-root operator.
Efficient variants of the WPM, including a recent differentiable formulation,
have been developed by ~\citet{Schmidt2016,BarreBrunel2025}. These
methods approximate the nonparaxial wave equation in Fourier space,
enabling the use of the FFT while avoiding the need to choose a Taylor-series
cutoff, a noted drawback of the FC-MS method ~\citet{2025ChenDyckPaper}.
The WPM is therefore a promising starting point for a practical and
more efficient wide-angle multislice method, but to our knowledge
this adaptation has not yet been systematically assessed in the electron
microscopy context.

This overview motivates the aim of this work: to re-evaluate high-angle
propagation methods looking for alternative methods that would be
more accurate than the F-MS without being to computer demanding and
that thus could be efficiently used in modern inverse EM methods.

- - - - - - sentence that can be used for improving somewhere the
text ? JLR (to improve) using representative thick specimens at high
electron energies, and demonstrate an alternative high-angle propagator
to the FC-MS method (ignoring backscattering). In particular, we compare
AS-MS and WPM against both the traditional F-MS method and a forward-ODE
reference solution of the full second-order time-independent wave
equation. We quantify accuracy and computational cost and determine
whether either method can provide an efficient correction to paraxial
multislice for representative crystalline TEM and STEM simulations.

- - - - - - - - - - - - - sentences forming the end of the short JLR
version without David's good historical review - - - - -

This approximation is highly efficient and accurate for small scattering
angles, but its limitations become increasingly important when high-angle
components contribute significantly to the exit wave, for example
in thick specimens, high-Z materials, or strongly scattering objects
or low energy EM ~\citet{RotherScheerschmidt2009,hillier_role_2017,WackerSchroeder2015,ming2013pcms}.

Since modern quantitative and inverse EM methods rely directly on
the fidelity of the propagation model, it is timely to revisit this
propagation step and assess whether a more accurate yet still efficient
alternative can improve the standard Fresnel multislice framework.
- -- - - - - - - - - - - end of sentences, kept just for thinking

In this work, we compare three multislice (MS) propagators: F-MS,
AS-MS, and WP-MS. Section~\ref{sec:algorithms_electron_matter} first
reviews the wave equation and the main propagation algorithms used
to model electron--matter interactions. Section~\ref{sec:results}
then benchmarks the methods and presents their behaviour in representative
CBED simulations. Finally, Section~\ref{sec:discussion} discusses
when each propagator method is expected to be adequate.

\section{Wave-Equation Models and Propagation Algorithms}

\label{sec:algorithms_electron_matter}

This section introduces the scalar wave-equation models and numerical
propagation algorithms considered in this work. We first define the
Klein--Gordon and relativistically corrected Schrödinger models (Subsection~\ref{sec:schrodinger_helmholtz}).
We then introduce F-MS and AS-MS, which differ only in their free-space
propagation kernel (Subsection~\ref{sec:fresnel_angular_spectrum}).
Finally, we discuss wide-angle and reference propagation methods,
including WP-MS and the ODE calculation used for benchmarking (Subsection~\ref{sec:advanced_propagation_methods}).

Throughout this work we use the convention now commonly adopted in
the electron microscopy literature, as used for example by Kirkland
and Chen--Van Dyck \citet{kirkland2010advanced,ChenVanDyck1997}.
$e$ will design the absolute value of the negative electron charge
as it is done in kirkland and Reimer; so $e$ is positive. But more
importantly, spatial frequencies are measured in cycles per unit length
rather than radians per unit length. This convention differs from
that used in some older electron microscopy papers, including the
original multislice formulation of Cowley and Moodie~\citet{cowley1957scattering},
in Wacker and Schröder \citet{WackerSchroeder2015}, and in much of
the optics literature \citet{goodman}. The vacuum spatial frequency
$p_{0}$ is therefore 
\begin{equation}
k_{0}=\frac{1}{\lambda},
\end{equation}
where $\lambda$ is the relativistic electron wavelength. This differs
from the optics convention $K_{0}=2\pi/\lambda$, so that $K_{0}=2\pi k_{0}$.

The consequence of this convention is twofold. First, we define the
Fourier transform with the phase factor 
\begin{equation}
\mathcal{F}\{f(x,y)\}(q_{x},q_{y})=\int f(x,y)\exp\left[-2\pi i(q_{x}x+q_{y}y)\right]\,dx\,dy,\label{eq:fourier_convention}
\end{equation}
where $(q_{x},q_{y})$ are transverse spatial frequencies in cycles
per unit length. Second, a plane wave propagating in the direction
$\mathbf{k}=(k_{x},k_{y},k_{z})$, with $|\mathbf{k}|=k_{0}=1/\lambda$,
is written as

\begin{equation}
\psi(\mathbf{r})=\exp\!\left(2\pi i\,\mathbf{k}\cdot\mathbf{r}\right)=\exp\!\left[2\pi i\left(k_{x}x+k_{y}y+k_{z}z\right)\right].\label{eq:plane_wave_convention}
\end{equation}


\subsection{Scalar Relativistic Wave Equations}

\label{sec:schrodinger_helmholtz}

As demonstrated by Fujiwara~\citet{fujiwara1961relativistic} and
detailed in Appendix~\ref{app:relativistic_hamiltonian}, the propagation
of high-energy electrons in matter can be described by the time-independent,
relativistically corrected Schrödinger equation 
\begin{equation}
\left[-\frac{\hbar^{2}}{2\gamma m_{0}}\nabla^{2}+E_{\mathrm{pot}}(\mathbf{r})\right]\psi(\mathbf{r})=\frac{p^{2}_{0}}{2\gamma m_{0}}\psi(\mathbf{r}),\label{eq:corrected_schrodinger_main}
\end{equation}

where $\psi(\mathbf{r})$ is the spatial part of the electron wave
function, $\hbar$ is the reduced Planck constant, $m_{0}$ is the
electron rest mass, $\gamma$ is the Lorentz factor, $p_{0}$ is the
electron momentum in vacuum, and $E_{\mathrm{pot}}$ is the potential
energy associated with the electrostatic potential $V$, expressed
in volts. With $e>0$ denoting the elementary charge,

\begin{equation}
E_{\mathrm{pot}}(\mathbf{r})=-eV(\mathbf{r}).\label{eq:potential_energy_main}
\end{equation}
Compared with the non-relativistic Schrödinger equation, the two relevant
changes are that the mass in the kinetic operator is the relativistic
mass $\gamma m_{0}$, and that the vacuum carrier wave is fixed by
relativistic energy and momentum relations: 
\begin{equation}
\begin{aligned}E_{0} & =m_{0}c^{2}, & E_{\mathrm{tot}} & =\gamma E_{0},\\
E_{\mathrm{kin}} & =E_{\mathrm{tot}}-E_{0}=(\gamma-1)E_{0}, & \gamma & =\frac{1}{\sqrt{1-v^{2}/c^{2}}},
\end{aligned}
\label{eq:relativistic_energy_relations}
\end{equation}
\begin{equation}
\begin{aligned}p_{0} & =\gamma m_{0}v=\frac{1}{c}\sqrt{E^{2}_{\mathrm{tot}}-E^{2}_{0}}=\frac{h}{\lambda}=hk_{0},\\
k_{0} & =\lambda^{-1}.
\end{aligned}
\label{eq:relativistic_momentum_relations}
\end{equation}

The right-hand side of Eq.~\eqref{eq:corrected_schrodinger_main},
$p^{2}_{0}/(2\gamma m_{0})$, is therefore the eigenvalue of the Schrödinger-like
stationary equation; it should not be confused with the total relativistic
energy $E_{\mathrm{tot}}$, which appears in the time factor $\exp(-iE_{\mathrm{tot}}t/\hbar)$.

Equation~\eqref{eq:corrected_schrodinger_main} can be rewritten
in Helmholtz form, 
\begin{equation}
\nabla^{2}\psi(\mathbf{r})+(2\pi k_{0})^{2}n^{2}(\mathbf{r})\psi(\mathbf{r})=0.\label{eq:helmholtz_form}
\end{equation}
This representation provides a direct link between electron microscopy
and the broader optics community, where propagation algorithms are
commonly formulated for Helmholtz equations. As explained in Appendix~\ref{app:relativistic_hamiltonian},
Eq.~\eqref{eq:helmholtz_form} may be used either with the more general
Klein--Gordon refractive index $n_{\mathrm{KG}}$ or with its linearised-potential
approximation $n_{\mathrm{S}}$. The latter accurately describes the
usual electron-microscopy regime, for which $|E_{\mathrm{pot}}(\mathbf{r})|\ll E_{\mathrm{tot}}$.

Introducing the conventional TEM interaction constant \citet{ChenVanDyck1997,kirkland2010advanced}
\begin{equation}
\sigma=\frac{2\pi\gamma m_{0}e}{h^{2}k_{0}}=\frac{2\pi\gamma m_{0}e\lambda}{h^{2}},\label{eq:interaction_constant}
\end{equation}
the refractive index used in the corrected Schrödinger regime is 
\begin{equation}
n^{2}_{\mathrm{S}}(\mathbf{r})=1-\frac{2\gamma m_{0}}{h^{2}k^{2}_{0}}E_{\mathrm{pot}}(\mathbf{r})=1+\frac{\sigma}{\pi k_{0}}V(\mathbf{r}).\label{eq:schrodinger_refractive_index}
\end{equation}
Equivalently, to first order in the potential, $n_{\mathrm{S}}(\mathbf{r})\simeq1+\sigma V(\mathbf{r})/(2\pi k_{0})$.

\subsection{Multislice Propagators}

\label{sec:fresnel_angular_spectrum}

To obtain conventional multislice from Eq.~\eqref{eq:helmholtz_form},
the wavefunction is written as a rapidly oscillating carrier multiplied
by a slowly varying envelope, 
\begin{equation}
\psi(x,y,z)=\exp(2\pi ik_{0}z)\tilde{\psi}(x,y,z).\label{eq:carrier_envelope}
\end{equation}
Neglecting $\partial^{2}\tilde{\psi}/\partial z^{2}$ gives the paraxial
equation 
\begin{equation}
\frac{\partial\tilde{\psi}}{\partial z}=\frac{i}{4\pi k_{0}}\nabla^{2}_{\perp}\tilde{\psi}+i\pi k_{0}\left[n^{2}(x,y,z)-1\right]\tilde{\psi}=\frac{i}{4\pi k_{0}}\nabla^{2}_{\perp}\tilde{\psi}+i\sigma\tilde{\psi}.\label{eq:paraxial_helmholtz_intro}
\end{equation}
For the corrected Schrödinger model, Eq.~\eqref{eq:schrodinger_refractive_index}
reduces the interaction term to $i\sigma V\tilde{\psi}$. For slice
$j$, define the projected electrostatic potential and transmission
function by 
\begin{equation}
V_{j}(x,y)=\int^{z_{j}+\Delta z}_{z_{j}}V(x,y,z)\,dz,\qquad t_{j}(x,y)=\exp\left[i\sigma V_{j}(x,y)\right].\label{eq:multislice_transmission}
\end{equation}

Fresnel multislice (F-MS) applies this transmission function and then
propagates the wave through vacuum with the paraxial Fresnel kernel.
Dropping the tilde, one slice is 
\begin{equation}
\psi_{j+1}(x,y)\approx\mathcal{F}^{-1}\left[\exp\left(-i\pi\frac{k^{2}_{\perp}}{k_{0}}\Delta z\right)\mathcal{F}\{t_{j}(x,y)\psi_{j}(x,y)\}\right].\label{eq:fresnel_step}
\end{equation}
Each slice requires one FFT pair. The paraxial approximation enters
through 
\[
k_{z}=\sqrt{k^{2}_{0}-k^{2}_{\perp}}\simeq k_{0}-\frac{k^{2}_{\perp}}{2k_{0}}.
\]
After removal of the carrier phase, this expansion gives the kernel
in Eq.~\eqref{eq:fresnel_step}. F-MS therefore combines a split-step
representation of the specimen with a small-angle approximation to
free-space propagation.

Angular-spectrum multislice (AS-MS) keeps the same transmission function
and split-step structure, but replaces the Fresnel kernel by the exact
free-space dispersion relation~\citet{goodman,ming2013pcms}: 
\begin{equation}
\psi_{j+1}(x,y)\approx\mathcal{F}^{-1}\left[\exp\left(2\pi i\left[\sqrt{k^{2}_{0}-k^{2}_{\perp}}-k_{0}\right]\Delta z\right)\mathcal{F}\{t_{j}(x,y)\psi_{j}(x,y)\}\right].\label{eq:asm_step}
\end{equation}
Apart from this Fourier-space phase, AS-MS is identical to F-MS and
retains the same $\mathcal{O}(P\log P)$ per-slice cost for $P$ transverse
pixels. It removes the paraxial free-space approximation, while retaining
the thin-slice split-step treatment of the specimen.

\subsection{Wide-Angle and Reference Propagation Methods}

\label{sec:advanced_propagation_methods} - - - - - - - - - - - -
to keep here or put in intro ??

F-MS is usually adequate for conventional high-energy TEM and STEM
simulations of thin, low-$Z$ specimens, where other uncertainties
such as thermal diffuse scattering, detector response, or microscope
aberration characterization often dominate the comparison with experiment.
However, for low incident electron energies, thick specimens, or high-$Z$
materials, F-MS can exhibit significant errors because of its paraxial
approximation and omission of backscattering ~\citet{RotherScheerschmidt2009,hillier_role_2017,WackerSchroeder2015,ming2013pcms}.
AS-MS is the minimal FFT-based correction to the paraxial propagation
step, but it still treats the specimen as a sequence of thin transmission
functions separated by vacuum propagation. More complete wide-angle
methods therefore seek to improve the propagation inside the specimen,
or to use the full second-order wave equation as a reference.

Further improvements in accuracy, particularly the treatment of large-angle
scattering and backward-propagating components, incur significant
computational cost. A more accurate electron microscopy multislice
algorithm is the Chen--Van Dyck method, also known as Fully Corrected
Multislice (FC-MS) \citet{ming2013pcms,ChenVanDyck1997}. The complexity
of FC-MS arises because the full second-order wave equation contains
a square-root operator that mixes a transverse derivative term with
the potential and cannot be evaluated as a single Fourier-space propagator.

FC-MS addresses this problem by expanding the square-root operator
in a Taylor series and retaining higher-order correction terms that
couple the potential to transverse derivatives. In this way the method
accounts more accurately for how the wave evolves within the potential
of each slice, rather than treating the slice only as a projected
phase grating followed by free-space propagation. This allows improved
treatment of larger scattering angles, low-energy electrons, thicker
specimens, and high-$Z$ materials and, in its more complete formulation,
also gives access to backward-propagating components.

5 years before the FC-MS was introduced in the field of scalar light
optics, a solution to the identical mathematical problem of solving
a Time Independent Wave Equation to higher accuracy than the F-MS
or AS-MS method had already been proposed, known as the Padé Beam-Propagation
Method (Padé BPM) ~\citet{Hadley1992}. In light optics, the time-independent
Helmholtz equation also leads to a square-root propagation operator
when it is factorized into forward- and backward-propagating branches.
Padé beam-propagation methods approximate this square-root operator
using rational functions rather than the Taylor polynomials proposed
by Chen and Van Dyck. Because a Padé approximant is a ratio of polynomials,
it can represent the square-root dispersion relation more accurately
over a wider angular range than a truncated Taylor series of the same
order. In practice, this often gives improved stability and accuracy
for large-angle propagation, allowing lower-order Padé approximants
to achieve accuracy that would otherwise require higher-order Taylor
corrections. In either case, approximating the square-root operator
remains computationally expensive and incurs significant overhead
relative to F-MS or AS-MS.

In electron optics, a still more accurate reference can be obtained
by directly integrating the full second-order time-independent wave
equation as an ordinary differential equation (ODE). Numerous authors
have used such solutions as references for assessing the accuracy
of F-MS and the FC-MS method. Wacker and Schröder provide a clear
derivation of the ODE and operator-approximation methods and benchmark
them against conventional F-MS and direct Runge--Kutta integration
of the second-order Schrödinger equation. Their results again showed
that the F-MS method can introduce substantial errors at low voltages,
large specimen thicknesses, and high scattering angles. Rother and
Scheerschmidt approached the problem from a relativistic perspective,
using ODE solutions of the Dirac, Klein--Gordon, and relativistically
corrected Schrödinger wave equations and comparing them with a paraxial
multislice solution. Their work similarly indicates that the dominant
high-angle error arises from the parabolic propagation step used in
F-MS. Hillier et al.~\citet{hillier_role_2017} show also that the
second order derivative term is important to track the positions of
Laue Zone rings in low energy CBED patterns of Silicon at 20keV. However,
out of all the methods discussed so far, it is only \citet{ming2013pcms}
that attempted to account for high angle effects in a way that can
exploit the Fourier transform, with their spherical propagation method.
The other series expansion and ODE methods cannot exploit the FFT
as directly as split-step multislice, limiting their practical use
as a forward model within inverse reconstruction algorithms.

The light-optics community has sought alternatives that retain the
computational efficiency of the FFT while avoiding the numerical difficulties
of evaluating Padé BPM approximants and still providing a nonparaxial
solution to the time-independent wave equation. One such solution
to this problem is the Wave Propagation Method (WPM), introduced by
Brenner in 1993~\citet{Brenner1993}. In the WPM, the potential or
refractive index medium is treated as locally homogeneous over small
transverse regions, and a corresponding homogeneous one-way propagator
is applied for each representative refractive index (see \ref{fig:wpm_explainer}
for a visual explanation of the method). The resulting fields are
then recombined according to the local refractive-index distribution.
This retains the potential-dependent wavelength within the propagation
step without explicitly forming the full dense square-root operator.
Efficient variants of the WPM, including a recent differentiable formulation,
have been developed by ~\citet{Schmidt2016,BarreBrunel2025}. These
methods approximate the nonparaxial wave equation in Fourier space,
enabling the use of the FFT while avoiding the need to choose a Taylor-series
cutoff, a noted drawback of the FC-MS method ~\citet{2025ChenDyckPaper}.
The WPM is therefore a promising starting point for a practical and
more efficient wide-angle multislice method, but to our knowledge
this adaptation has not yet been systematically assessed in the electron
microscopy context.

- - - - - -- -- - -- -

\subsubsection{Classical Wave Propagation Method in Scalar Optics}

\label{sec:wpm_optics}

\begin{figure}[H]
\centering \safeincludegraphics[width=0.95\textwidth,keepaspectratio]{method_explainer_generations/wpm}
\caption{Schematic of the Wave Propagation Method (WPM) as presented by \citet{Schmidt2016}.
The input wave is propagated using a bank of homogeneous angular-spectrum
kernels, each evaluated at a different representative refractive index
$n_{m}$. At each output pixel, the propagated field corresponding
to the local refractive index is selected via mask and recombined.
Our improvement to this algorithm (not described in this figure for
simplicity) includes a soft binning step which makes this method suitable
for a continuous refractive index distribution - see Figure \ref{fig:wpm_binning_diagnostics}
below.}
\label{fig:wpm_explainer} 
\end{figure}

The Wave Propagation Method (WPM) gathers the split-step approximations
in F-MS AS-MS in a one-step propagation, by including the phase shift
of the potential in the propagation kernel itself. Indeed, as it is
explained below, starting from the local refractive index of the medium,
new local wave vectors $\mathbf{k}(x,y)$ are defined and use in the
propagation step.

Returning to the full second-order Helmholtz equation \eqref{eq:helmholtz_form},
it can be rearranged into an operator equation with a forward and
backward branch \citet{Schmidt2016}. Ignoring the backward branch,
the forward branch may be written as 
\begin{equation}
\partial_{z}\psi=2\pi i\left[\sqrt{k^{2}_{0}n^{2}(x,y,z)+\frac{1}{4\pi^{2}}\nabla^{2}_{\perp}}-k_{0}\right]\psi.\label{eq:one_way_square_root_intro}
\end{equation}
The central concern is that the square-root operator inside the wave
equation mixes the potential and the transverse derivative, and cannot
be evaluated as a single Fourier-space propagator. Padé BPM and FC-MS
methods solve for this operator using rational or series approximations
to the square-root term, which are expensive and cannot usually directly
use the FFT.

Looking again at the AS-MS propagation kernel of equation \eqref{eq:asm_step},
\[
\left[\sqrt{k^{2}_{0}-k^{2}_{\perp}}-k_{0}\right]
\]
what is primarily missing is a refractive-index term. We can reintroduce
this term locally by replacing $k_{0}$ with $k_{0}n(x,y)$, giving

\[
\left[\sqrt{k^{2}_{0}n^{2}(x,y)-k^{2}_{\perp}}-k_{0}\right].
\]

For a single refractive-index value, this kernel can be evaluated
straightforwardly with the angular spectrum method. The WPM simply
extends the angular spectrum method by evaluating a bank of homogeneous
angular-spectrum propagators for each refractive index in the slice,
and then retaining the propagated field corresponding to the local
refractive index at each output pixel.

This operation is obviously simplistic, because output pixels can
receive a propagated field that travelled from some pixels with the
incorrect refractive index. However, the approximation lives in the
fact that those different refractive indices are further away from
the local refractive-index value of that pixel, and therefore contribute
less to the final result. In this manner, the propagation step accounts
for the local wavelength inside the specimen without explicitly evaluating
the full square-root operator. Schmidt et al.~\citet{Schmidt2016}
showed that the leading error from neglecting coupling between neighbouring
refractive-index regions affects amplitude rather than phase; consequently,
the WPM phase propagation is correct up to second order.

Figure~\ref{fig:wpm_explainer} shows the same local-homogeneous
idea in the form used by our electron-microscopy adaptation below.
In the classical WPM as implemented by Schmidt, a mask is created
to separate out each unique refractive index in a slice, the angular
spectrum kernel is calculated for each refractive index, and then
the mask is used to choose the correct angular spectrum solution from
the bank of all computed propagations.

In Brenner's original formulation~\citet{Brenner1993}, evaluating
a different propagation kernel at every output pixel is prohibitively
expensive. For a transverse grid with $P$ pixels, the direct mixed
real/Fourier-space calculation scales as $\mathcal{O}(P^{2})$. Schmidt
et al.~\citet{Schmidt2016} made the method practical for many optical
simulations by grouping pixels with the same refractive index. Only
one homogeneous angular-spectrum propagation is then required for
each distinct index value. If a slice contains $M$ such values, the
cost is reduced to approximately 
\[
\mathcal{O}(MP\log P),
\]
instead of a direct all-to-all evaluation.

This saving is large in light-optical simulations where a slice may
contain only a small number of unique refractive indices. It is less
immediate in electron microscopy, where the electrostatic potential
varies continuously and almost every pixel may have a different refractive
index. Exact grouping therefore gives little benefit unless the refractive-index
field is deliberately discretised. A second issue is that hard grouping
introduces sharp masks, which are inconvenient for gradient-based
optimisation because the assignment of pixels to refractive-index
groups is not smoothly differentiable.

\subsubsection{Wave Propagation Multislice}

\label{sec:wp_ms}

We refer to our adaptation of the WPM suitable for electron microscopy
as wave propagation multislice (WP-MS), to distinguish it from the
classical WPM used in scalar light optics. To make the computation
of WP-MS tractable in the context of electron microscopy, and to avoid
the naive $O(N^{2D})$ calculation for a $D$-dimensional transverse
grid, we adapt partially the WPM algorithm presented in~\citet{BarreBrunel2025}.
Their method makes use of a soft-partition mask between a slice with
two unique refractive index values to maintain differentiability of
the WPM algorithm. However, they note that they can only use this
mask because they have chosen an object with two refractive index
values, and an extension to a more general case would be more challenging.

To adapt this method to a slice with more refractive index values,
we first discretise the refractive-index range $[n_{\mathrm{min}},n_{\mathrm{max}}]$
of the current slice into $N_{\mathrm{bins}}$ discrete reference
values, denoted as $\{n^{(k)}_{\mathrm{ref}}\}^{N_{\mathrm{bins}}}_{k=1}$.

In a typical multislice potential, many of the refractive-index values
will be very close to $1$ in the vacuum region, while the refractive-index
values inside an atom will concentrate in a narrow range slightly
above $1$. Since there are likely to be many more refractive-index
values near the vacuum region, and the high refractive index values
of the atoms are likely to be more sparsely distributed (as the atoms
take up comparatively less space in a slice), we employ a quadratic
distribution for the reference indices rather than using linearly
spaced bins, in order to allocate more bins near $n=1$. The reference
indices are distributed according to a quadratic function determined
by the parameter $\gamma$, with $\gamma=2.0$ for quadratic spacing:

\begin{equation}
n^{(k)}_{\mathrm{ref}}=n_{\mathrm{min}}+(n_{\mathrm{max}}-n_{\mathrm{min}})\left(\frac{k-1}{N_{\mathrm{bins}}-1}\right)^{\gamma},\quad k=1,\ldots,N_{\mathrm{bins}},\label{eq:wpm_bin_spacing}
\end{equation}

which allows for denser binning near the vacuum refractive index when
$\gamma>1$ See \ref{fig:wpm_binning_diagnostics} for a visual diagram
showing graphically the binning step of a slice.

Next, for a propagation step $\Delta z$, we compute a bank of propagated
fields $u_{k}(x,y)$. Each field $u_{k}$ corresponds to the diffraction
of the input field $u_{\mathrm{in}}(x,y)$ assuming a homogeneous
medium with refractive index $n^{(k)}_{\mathrm{ref}}$. This is performed
efficiently using batch FFT operations with the angular spectrum kernel:

\begin{equation}
u_{k}(x,y)=\mathcal{F}^{-1}\left[\mathcal{F}\{u_{\mathrm{in}}\}\exp\left(2\pi i\Delta z\left[\sqrt{(k_{0}n^{(k)}_{\mathrm{ref}})^{2}-k^{2}_{x}-k^{2}_{y}}-k_{0}\right]\right)\right].\label{eq:wpm_binned_fields}
\end{equation}

To evaluate the field at a specific spatial coordinate $(x,y)$ where
the local refractive index is $n(x,y)$, we identify the two bounding
reference bins such that $n_{L}\le n(x,y)<n_{R}$, where $n_{L},n_{R}\in\{n^{(k)}_{\mathrm{ref}}\}$.
We adopt a differentiable interpolation scheme by defining a local
interpolation weight

\begin{equation}
w_{\mathrm{raw}}=\frac{n(x,y)-n_{L}}{n_{R}-n_{L}},\quad w_{\mathrm{raw}}\in[0,1]\label{eq:wpm_raw_weight}
\end{equation}

To ensure smooth gradients and minimise discretisation artifacts,
we map this linear weight through a smoothstep activation function
$S(w)$, akin to the spatial partitioning technique demonstrated by~\citet{BarreBrunel2025}
via a cubic Hermite interpolation function:

\begin{equation}
w=S(w_{\mathrm{raw}})=3w^{2}_{\mathrm{raw}}-2w^{3}_{\mathrm{raw}}\label{eq:wpm_smoothstep}
\end{equation}

While linear interpolation would satisfy the requirement for differentiability,
it could result in discontinuities in the gradient at the bin boundaries
in a gradient-based optimisation scheme.

The final field at the next step $u(x,y,z+\Delta z)$ is computed
as the weighted superposition of the fields propagated at the bounding
indices:

\begin{equation}
u(x,y)=(1-w)\cdot u_{L}(x,y)+w\cdot u_{R}(x,y)\label{eq:wpm_two_bin_interpolation}
\end{equation}

Equivalently, the WP-MS update can be written as 
\begin{equation}
\psi(x,y,z+\Delta z)\approx\sum^{N_{\mathrm{bins}}}_{m=1}w_{m}(x,y)\mathcal{F}^{-1}\left[\exp\left(2\pi i\left[\sqrt{k^{2}_{0}\left(n^{(m)}_{\mathrm{ref}}\right)^{2}-k^{2}_{\perp}}-k_{0}\right]\Delta z\right)\mathcal{F}\{\psi(x,y,z)\}\right].\label{eq:wpm_bank}
\end{equation}
This gives a differentiable approximation to the one-way medium propagator
in equation~\eqref{eq:one_way_square_root_intro}. When the slice
is homogeneous, only one refractive-index bin is needed and WP-MS
reduces to the angular spectrum method. For an inhomogeneous specimen,
the method is more expensive than AS-MS because many FFT-based propagations
are required per slice, but it accounts for the local wavelength inside
the specimen rather than propagating every slice with the vacuum wavelength.

This formulation ensures that the propagation model is fully differentiable
with respect to the refractive-index map $n(x,y)$, enabling inverse
reconstruction of potentials via automatic differentiation. It also
means that only a limited number of FFTs, equal to $N_{\mathrm{bins}}$,
are required per propagation step, significantly reducing the computational
cost compared to a naive WPM implementation. The accuracy of WP-MS
depends on the number of bins $N_{\mathrm{bins}}$ chosen, and will
change on a case-by-case basis depending on the refractive-index distribution
of the sample. Figure~\ref{fig:wpm_binning_diagnostics} shows a
representative diagnostic of the binning procedure. A practical way
to mitigate this potential source of error is to perform a test simulation
and double the number of bins until the required tolerance is achieved
in the slice error.

\begin{figure}[H]
\centering \safeincludegraphics[width=\textwidth,keepaspectratio]{wpm_binning_diagnostics}
\caption{Diagnostics of the differentiable WP-MS binning strategy for a representative
MoS$_{2}$ slice. Top row: discretized refractive-index map (bin lower
bound $n_{L}$), smooth interpolation weights $w$, and histogram
of the true refractive-index distribution with superimposed bin centers.
Bottom row: spatial approximation error $n_{\mathrm{true}}-n_{\mathrm{fit}}$,
error histogram on a log scale, and a 1D cross-section comparing the
original and fitted refractive-index profiles.}
\label{fig:wpm_binning_diagnostics} 
\end{figure}

In general, WP-MS will need a GPU for efficient computation, and we
have found that $N_{\mathrm{bins}}$ between 16 and 128 are sufficient
for the simulations presented in this work. This is a marked improvement
over the naive case, as it enables WP-MS to simulate many slices with
larger grid sizes ($512\times512$ or $1024\times1024$ pixels) in
a reasonable time frame (seconds to minutes), provided the refractive-index
binning strategy is capable of correctly capturing the refractive-index
distribution of the sample.

\subsubsection{Forward ODE Reference}

\label{sec:forward_ode_reference}

We use a direct Ordinary Differential Equation (ODE) solution of the
second-order Klein--Gordon Helmholtz equation as the reference calculation.
The ODE solution is a standard way to solve a wave equation numerically,
whereby the transverse coordinates are discretised, the second-order
equation is rewritten as a first-order system, and the resulting large
system of ordinary differential equations is integrated along $z$.
Substituting $n_{\mathrm{KG}}$ into the Helmholtz equation gives
\begin{equation}
\partial^{2}_{z}\psi=-\left[\nabla^{2}_{\perp}+(2\pi k_{0})^{2}n^{2}_{\mathrm{KG}}(x,y,z)\right]\psi.\label{eq:kg_second_order}
\end{equation}
Introducing $\phi=\partial_{z}\psi$ gives the equivalent first-order
system 
\begin{equation}
\frac{d}{dz}\begin{pmatrix}\psi\\
\phi
\end{pmatrix}=\begin{pmatrix}\phi\\
-\left[\nabla^{2}_{\perp}+(2\pi k_{0})^{2}n^{2}_{\mathrm{KG}}(x,y,z)\right]\psi
\end{pmatrix}.\label{eq:kg_state}
\end{equation}
We integrate this system with an adaptive high-order Runge--Kutta
method. The transverse Laplacian is evaluated spectrally using FFTs,
while the potential term is applied by pointwise multiplication in
real space.

The calculation is initialised on the forward-propagating vacuum branch.
The entrance wave $\psi(z=0)$ is supplemented by the longitudinal
derivative 
\begin{equation}
\phi(z=0)=\mathcal{F}^{-1}\left[2\pi i\sqrt{k^{2}_{0}-k^{2}_{\perp}}\,\mathcal{F}\{\psi(z=0)\}\right],\label{eq:kg_phi0}
\end{equation}
which selects the positive-$k_{z}$ branch of the incident field in
vacuum. We therefore refer to this as the ODE reference. The method
is much more expensive than the split-step propagation methods, although
it does include back-scattered effects that occur throughout the interaction
process.

The four methods we have presented thus far can be read as successive
changes to the same underlying Helmholtz problem, with a chosen refractive
index representing either the Klein--Gordon or the relativistically
corrected version of the time-independent wave equation, and then
an appropriate approximation to the solution of the wave equation
F-MS uses a split-step specimen approximation together with the paraxial
free-space propagator. AS-MS keeps the split-step approximation but
restores the exact free-space spherical propagator. WP-MS approximates
the one-way propagation operator inside the medium by local homogeneous
propagation. The ODE solution keeps the second-order Klein--Gordon
Helmholtz equation and evolves the transverse Laplacian and specimen
potential together, and it is used as the reference calculation in
the first numerical experiment. See Table \ref{tab:method_timeline}
below for a summary timeline of the methods presented in this work.

\begin{table}[htbp]
\centering {\small\setlength{\tabcolsep}{5pt} 
\global\long\def\arraystretch{1.15}%
 }{\small{}%
\begin{tabular}{@{}L{0.18\textwidth}L{0.25\textwidth}L{0.25\textwidth}L{0.22\textwidth}@{}}
\toprule 
{\small\textbf{Method family}}{\small{} } & {\small\textbf{Light-optics reference}}{\small{} } & {\small\textbf{Electron-optics reference}}{\small{} } & {\small\textbf{Description}}{\small{} }\tabularnewline
\midrule 
{\small Fresnel MS } & {\small Fresnel; Fleck--Morris--Feit \citet{fresnel1818memoire,fleck1976atmosphere} } & {\small Cowley--Moodie; Ishizuka--Uyeda \citet{cowley1957scattering,ishizuka1977practical} } & {\small Split-step propagation with the paraxial Fresnel kernel. }\tabularnewline
{\small Angular Spectrum MS } & {\small Booker--Clemmow; Korpel et al. \citet{booker1950angular,korpel1986splitstep} } & {\small Ming--Chen \citet{ming2013pcms} } & {\small Same split-step structure, but with the exact free-space propagator. }\tabularnewline
{\small Padé/Series Expansion (BPMs) } & {\small Hadley; Lee--Voges \citet{Hadley1992,lee1994wideangle} } & {\small van Dyck; Chen--Van Dyck \citet{vandyck1979higherorder,ChenVanDyck1997} } & {\small Higher-order approximations to the square-root propagation
operator. }\tabularnewline
{\small Wave Propagation Method } & {\small Brenner--Singer; Schmidt et al. \citet{Brenner1993,Schmidt2016} } & {\small This work } & {\small Local homogeneous angular-spectrum kernels applied inside the
medium. }\tabularnewline
{\small ODE } & {\small Banerjee--Sharma \citet{banerjee1989helmholtz,sharma1989totalfields} } & {\small Wacker--Schröder; Rother--Scheerschmidt\textsuperscript{{*}}
\citet{WackerSchroeder2015,RotherScheerschmidt2009} } & {\small Direct and reciprocal-space references for non-paraxial propagation. }\tabularnewline
\midrule 
\multicolumn{4}{@{}L{0.90\textwidth}@{}}{{\footnotesize\textsuperscript{{*}}Rother--Scheerschmidt used a
forward-scattering approximation in Fourier space for a periodic crystal,
rather than direct real-space integration of the full second-order
ODE.}}\tabularnewline
\bottomrule
\end{tabular}}{\small\caption{Historical timeline displaying correspondence between the propagation
methods used here and closely related developments in light and electron
optics}
\label{tab:method_timeline} }
\end{table}


\section{Benchmarks and Results}

\label{sec:results}

We present three numerical experiments designed to identify when each
propagation method is useful. The first reproduces the Au~{[}100{]}
beam-amplitude benchmark of Rother and Scheerschmidt~\citet{RotherScheerschmidt2009},
where F-MS, AS-MS, and WP-MS can be compared with an ODE solution
of the full second-order Klein--Gordon equation. After this validation,
the ODE reference is not used in the remaining tests because of its
high computational cost. The second experiment compares the FFT-based
methods for a lower-$Z$ Si~{[}111{]} CBED calculation at high energy
and large thickness. The third experiment considers a thicker Au~{[}100{]}
convergent-probe calculation, using WP-MS as the practical high-angle
reference.

\subsection{Au~{[}100{]} ODE Benchmark}

\label{sec:verification}

\begin{figure}[H]
\centering \safeincludegraphics[width=\textwidth,keepaspectratio]{Au_exit_wave_amplitudes}
\caption{The left columns show the real-space exit-wave intensity and phase
for the Au~{[}100{]} verification case. The rightmost column shows
the HRTEM image intensity after applying a representative aberration-corrected
300~keV transfer function with residual $C_{s}=10~\mu\mathrm{m}$,
defocus of $+6.0$~nm, a 12~mrad objective aperture. Color scales
are shared across methods within each column, and image-plane NRMSE
values are relative to the Klein--Gordon ODE reference.}
\label{fig:au_exit_wave_amplitudes} 
\end{figure}

To verify our implementation, and to assess which propagation methods
retain accuracy at high angles in a thick high-$Z$ specimen, we recreated
the simulation presented in figure~3 of \citet{RotherScheerschmidt2009}.
This benchmark propagates a plane wave through an Au~{[}100{]} crystal
at 300~keV and compares against a converged ODE solution of the full
second-order Klein--Gordon equation. Periodic lateral boundary conditions
are implicit because each method uses Fourier propagation, so the
crystal is repeated in the transverse directions. To test thicker
objects, we extended the Rother--Scheerschmidt simulation from 25
unit cells ($\sim$10~nm) to 100 unit cells ($\sim$40~nm).

Figure~\ref{fig:au_exit_wave_amplitudes} shows the exit-wave intensity,
phase, and transfer-function image intensity after propagation through
100 unit cells of Au~{[}100{]}. The F-MS exit-wave amplitude and
phase are qualitatively different from the other three solutions,
whereas AS-MS, WP-MS, and the Klein--Gordon ODE reference remain
visually similar. After a representative microscope contrast transfer
function is applied, however, the differences between methods are
much less obvious. This explains why conventional F-MS can still be
adequate for qualitative comparison of atomic-column images, even
when the unfiltered exit wave contains appreciable high-angle error.
To make the comparison more quantitative, Figure~\ref{fig:au_pixel_amplitudes}
plots the normalised Fourier amplitude after each slice for the central
{[}0, 0{]} component and for the high-angle {[}0, 28{]} component
at approximately 135~mrad as was demonstrated in the work of \citet{RotherScheerschmidt2009}.

\begin{figure}[H]
\centering \safeincludegraphics[width=\textwidth,keepaspectratio]{Au_beam_amplitudes_full}
\caption{Pixel-amplitude comparison for the Au~{[}100{]} verification case.
The central and high-angle pixels are compared across F-MS, angular
spectrum multislice, wave propagation multislice, and the Klein--Gordon
ODE reference.}
\label{fig:au_pixel_amplitudes} 
\end{figure}

For the central {[}0, 0{]} component, all methods agree closely up
to approximately 10~nm, or 25 unit cells. Beyond this thickness,
the F-MS result begins to deviate from the other methods even on the
on-axis pixel. By about 15~nm, high-angle diffracted beams have propagated
far enough through the crystal to now interfere back into the central
beam, so an error in their accumulated phase can also modify the on-axis
amplitude. For the high-angle {[}0, 28{]} component at 135~mrad,
the F-MS solution diverges from the ODE reference after approximately
10 unit cells, or 4~nm, consistent with the findings of \citet{RotherScheerschmidt2009},
who attributed the discrepancy to the linearisation of the square-root
operator. AS-MS tracks the ODE reference much more closely than F-MS
and remains comparable to WP-MS. The WP-MS solution is slightly closer
to the reference, but the difference from AS-MS only becomes clear
at the largest thicknesses considered here. These results verify that
both the AS-MS method and WP-MS method can adequately track a complete
ODE solution to the wave equation, and provide an appreciable improvement
over the F-MS method, without the need to solve a higher order taylor
expanded method like the FC-MS solution presented by \citet{ChenVanDyck1997}.

\subsection{Si~{[}111{]} and Au~{[}100{]} CBED Simulations}

\label{sec:silicon_cbed}

After verifying AS-MS and WP-MS against the ODE solution, we focus
on the three FFT-based multislice methods: F-MS, AS-MS, and WP-MS.
The aim is to compare how their predictions differ in CBED patterns
for representative low-$Z$ and high-$Z$ specimens. We first simulate
Si~{[}111{]} at 100~keV with an 8~mrad probe and a specimen thickness
near 1000~\AA, following the spirit of the low-energy CBED comparison
of Hillier et al.~\citet{hillier_role_2017} but at a higher accelerating
voltage. The maximum scattering angle was set to 300~mrad.

\begin{figure}[H]
\centering \safeincludegraphics[width=\textwidth,keepaspectratio]{cbed_kirkland_si111_combined_300mrad_999A}
\caption{Si~{[}111{]} CBED comparison at 100~keV, 8~mrad probe, and approximately
1000~\AA{} thickness. (a) WP-MS CBED pattern cropped to 300~mrad.
(b) Radially integrated annular intensity for all three methods, with
detected diffraction peaks marked. (c) Peak displacement relative
to WP-MS, with a scaled parabolic-versus-exact dispersion error curve
overlaid. The F-MS---WP-MS shift (blue) follows the dispersion-error
curve closely, while the AS-MS---WP-MS shift (orange) remains near
zero, consistent with the Au~{[}100{]} results.}
\label{fig:si111_cbed_combined} 
\end{figure}

Figure~\ref{fig:si111_cbed_combined} shows the WP-MS CBED pattern,
the radial profiles with detected peaks, and the peak displacements.
The F-MS peaks shift progressively with scattering angle, but the
AS-MS tracks WP-MS closely, with peak displacements near zero across
the full 300~mrad range. This suggests that the dominant error in
F-MS is the free-space dispersion approximation rather than the split-step
treatment of the specimen, even for a lower-$Z$ material at a moderate
thickness. \citet{hillier_role_2017} have already identified the
need to include the second order derivative of the wave equation to
track the Laue zone positions to high angles, and explained using
the ewald sphere why there is a difference between traditional F-MS
and a second order ODE solution. Our results in part b demonstrate
the same effect, but without the need to solve an expensive and slow
second order ODE. In part c of \ref{fig:si111_cbed_combined}, we
show the x shift of each detected Laue zone peak, overlayed with a
rescaled dispersion relation phase error - i.e the difference between
a parabola and a sphere. The difference in peaks is proportional to
this phase error. The dispersion phase error is computed as $|\Delta\phi|=(2\pi/\lambda)\,\bigl|\sqrt{1-\lambda^{2}q^{2}}-1+\tfrac{1}{2}\lambda^{2}q^{2}\bigr|$,
where $q=\sin\theta/\lambda$ is the transverse spatial frequency:
the first term is the exact spherical dispersion $\sqrt{1-x}$ with
$x=\lambda^{2}q^{2}$, and the remaining terms are its first-order
Taylor expansion $1-x/2$, i.e.\ the Fresnel kernel. The curve in
panel~(c) is this expression fitted to the F-MS $-$ WP-MS peak shifts.
The shifted peaks and the fitted curve overlap well, indicating that
the error in the output is due to not tracking the spherical propagation
between slices in free space. The extra accuracy the WP-MS method
provides by tracking the change in wavelength through the potential
is not noticeable, and simply swapping the F-MS propagator to the
AS-MS propagator already reclaims much of the lost accuracy. It is
important to note that in general, thermal diffuse scattering will
significantly dampen the presence of peaks beyond 200 mrad, and therefore
the application of these methods will be limited for now in traditional
TEM, unless perhaps the sample is cooled.

Next, in order to test if a sample with a higher atomic number would
highlight more differences between each of the three methods that
can use the Fourier Transform, we simulated convergent-beam electron
diffraction from an Au~{[}100{]} supercell at 200~keV using a 5~mrad
aberration-free probe. The probe was propagated to a specimen thicknesses
of approximately 100~nm again. To sample the high-angle components
of the convergent beam accurately, the enlarged $9\times9$ gold cell
was discretized on a $2048\times2048$ grid. We then propagated the
wave using F-MS, AS-MS, and WP-MS, and compared the exit-surface CBED
patterns obtained from the Fourier transform of each exit wave.

\begin{figure}[H]
\centering \safeincludegraphics[width=\textwidth,keepaspectratio]{cbed_combined_amplitude_log_error_maps}
\caption{CBED exit-wave amplitude comparison for the Au~{[}100{]} convergence-angle
experiment. The maps show the final diffraction patterns produced
by F-MS, AS-MS, and WP-MS, together with log-scale amplitude errors
relative to WP-MS.}
\label{fig:cbed_combined_amplitude_log_error_maps} 
\end{figure}

Figure~\ref{fig:cbed_combined_amplitude_log_error_maps} shows the
same pattern in the recorded diffraction plane. The F-MS CBED amplitude
differs visibly from WP-MS, especially in the higher-angle features
of the CBED pattern. AS-MS remains much closer to WP-MS, and the remaining
difference in amplitude between AS-MS and WP-MS is small compared
with the discrepancy introduced by the Fresnel propagator.

\begin{figure}[H]
\centering \safeincludegraphics[width=\textwidth,keepaspectratio]{cbed_error_vs_thickness_au100}
\caption{Relative CBED amplitude error as a function of specimen thickness
for the Au~{[}100{]} convergence-angle study. The error is reported
against the WP-MS reference at the end of each unit cell.}
\label{fig:cbed_error_vs_thickness} 
\end{figure}

In Figure~\ref{fig:cbed_error_vs_thickness} we visualise the root
mean square (RMSE) error of the CBED amplitude against the WP-MS reference
as a function of thickness. After a thickness of approximately 10~nm,
there is an order-of-magnitude error difference between the F-MS and
AS-MS solutions. The AS-MS method does become less accurate as the
thickness increases, but remains much closer to WP-MS than to F-MS
throughout the calculation. The residual amplitude difference between
AS-MS and WP-MS can be explained as the WP-MS propagation kernel uses
the local refractive index, and therefore tracks the potential-dependent
wavelength through each slice, decreasing the final error in amplitude.
The refractive index of the gold atom near the centre is approximately
1.07, large enough such that over a sample of this thickness, tracking
this change can start to materially effect the output CBED image.
A similar residual difference in amplitudes between the WP-MS and
AS-MS is visible in the exit-wave comparison of Figure~\ref{fig:au_exit_wave_amplitudes}
from the first numerical experiment, even at 40nm.

Although ptychographic reconstruction for specimens as thick as this
is not yet feasible, our results show that the choice of propagator
can change the diffraction-plane amplitudes. Because ptychographic
reconstructions are sensitive to systematic amplitude errors, such
errors could be absorbed as artifacts or bias the recovered potential.
Therefore the choice of propagator may become more consequential as
ptychography is pushed toward thicker specimens and higher scattering
angles. Given that there is almost no cost associated with the AS-MS
method, we at least recommend that it become the default propagator
in almost all multislice use cases going forward to minimise another
possible source of error in the forward model.

The practical advantage of AS-MS is only immediately useful if it
does not require a much finer grid than conventional F-MS. The additional
sampling concern comes from the angular-spectrum propagation phase,
which is steeper than the Fresnel phase at large transverse wavevector.
For one propagation step,

\[
\phi_{\mathrm{F}}(k_{\perp})=-\pi\frac{\Delta z}{k_{0}}k^{2}_{\perp},\qquad\phi_{\mathrm{ASM}}(k_{\perp})=2\pi\Delta z\left[\sqrt{k^{2}_{0}-k^{2}_{\perp}}-k_{0}\right].
\]

Comparing the phase gradients gives the relative axial sampling penalty

\[
\frac{\left|d\phi_{\mathrm{ASM}}/dk_{\perp}\right|}{\left|d\phi_{\mathrm{F}}/dk_{\perp}\right|}=\frac{k_{0}}{\sqrt{k^{2}_{0}-k^{2}_{\perp}}}.
\]

This factor is only large when $k_{\perp}$ approaches $k_{0}$, where
the propagating spectrum approaches the evanescent boundary. It can
be written directly in terms of scattering angle because $k_{\perp}/k_{0}=\sin\theta$:

\[
\frac{k_{0}}{\sqrt{k^{2}_{0}-k^{2}_{\perp}}}=\frac{1}{\sqrt{1-\sin^{2}\theta}}=\frac{1}{\cos\theta}.
\]

Under the TEM conditions used here this remains close to unity. For
a 200~mrad probe, the largest correction to the Fresnel slice-thickness
condition is

\[
\frac{1}{\cos(200~\mathrm{mrad})}\approx1.02,
\]

so the axial sampling bound is tightened by only about 2\%. The transverse
grid sets a separate limit as it determines the largest scattering
angle that can be represented without aliasing. A real-space pixel
size $\Delta x$ gives a maximum transverse spatial frequency $k_{\perp,\mathrm{Nyq}}=1/(2\Delta x)$.
Since $k_{\perp}=k_{0}\sin\theta$ and $k_{0}=1/\lambda$, the corresponding
Nyquist angle is

\[
\theta_{\mathrm{Nyq}}=\arcsin\!\left(\frac{\lambda}{2\Delta x}\right)\approx320~\mathrm{mrad}.
\]

Here we used $\lambda\approx0.025\,\text{\AA}$ at 200~keV and $\Delta x=0.04\,\text{\AA}$.
This angle is still far from the evanescent boundary at $\theta=90^{\circ}$,
where the square root in the ASM kernel becomes singular. (Jean Luc
needs to add a section here about sampling)

\section{Discussion}

\label{sec:discussion}

The discussion part will adress several points. Firstly the interest
of the ODE will be discussed. Secondly, the surprising accuracy of
WP-MS will be analysed. Thirdly, the comparaison between F-MS and
AS-MS will be detailed. The conclusion part will summarize of the
important consequences of these discussions and results.

\subsection{ODE Discussion}

\label{sec:ODEdiscussion}

ODE is the best numerical solution to the generalized static Helmoz
equation, but in fact there exists many ways to solve an ODE equations.
This is a central problem physics and engineering and powerful mathemiatical
and numerical tools have been developp to solve them. In engineering,
finite element calculations is certainly the better instrument to
solve these kind of ODE. For instance ODE equations are at the core
the comsol project and language paradism like python or matlab have
develop similar FE solutions. One issue of this approach it is that
they are very technical and many physicists and electron microscopists
are not specialist of these fields. PROS : black boxes (COMSOL) are
available to solve this type of problem (pourquoi ne pas le faire
justement avec les licenses du LETI et via Matthieu pour un autre
papier ???? ) CONS : c1) mesh size have to be small to retain all
the accuracy c2) consequently if one wants to have a resolution approaching
$\lambda$, which will be designated as $\lambda$-simulations, hudge
meshed have to be used and computation time is enormous. c3) it seems
difficult to use these approaches even in simple standard $\lambda$-simulations
of perfect crystals. So using this kind of simulations in phase-retrieval
algorithms with a resolution approaching $\lambda$, like ptychography
is presently unpracticable. This is why more numerical efficient algorithm
like F-MS have to be developped. (6h54 est-ce que je continue d'écrire
les idées sur cette partie ODE ? non je crois que j'ai l'essentiel
pour l'instant donc je poursuis à la parite suivante NON j'oubliais
l'autre idée ) c4) Another annoying point is that most of the publications
citing ODE resolution of the Helmoz equation does not detail in length
the ODE solver that have and it is then difficult to reproduce their
result. This is why in this paper we have choosen to give all the
details of our ODE solver in the supplementary github site c5) In
this paper the ODE solver was used only to validate the 3 other methods
on a very simple gold crystal and see the main approximations of the
3 other methods. As ODE results is very close, to other methods even
for this high-Z we think it does not need to repeat the ODE for the
other systems that are study in this paper.

\subsection{WP-MS approach}

\label{sec:WP-MSapproach}

- surprising how it works well in spite of it crude approach ... -
explanations could be that (i) it can work well in EM because the
index n in matter is rather constant every except at the core the
atoms and (ii) it turns out that this approximation is not worse,
but even apparently better, that the traditionnal MS. Supposing the
wave length in matter is constant is also a very crude and simplistic
supposition, but it turns to be working quite well in many situations
of EM (see next sub-chapter) - the core of physics is to develop model
that mimic the results of experiments. Most of the time, these models
are a simplification of what really happen in matter but the interest
of the model is not the representation of matter but their habiities
to reproduce matter behavior for a given experiment. Models approach
the behavior of matter in certain conditions but are not a perfect
complete description of matter. A model becomes interesting when it
is efficient and simple. - a point interesting to discuss is the change
of wavelength that happen in the WP-MS as it also happen in Bloch
wave methods. Is that in change in wavelength physical ? Does it really
happen in matter ? If yes, how coherent, a key point in ptychography,
can be preserved ? This paper will not give an exhausting answer to
this question. We will just give the starting following analysis :
the change of a wavelength in WP-MS is the just a technical and simple
representation of what happens in the experiment and it can be said
that in this static representaiton of beam propagation there is indeed
a change of wavelength in the piece of sample. But this wave length
is not really the wavelength that is linked to coherence : coherence
and interference effect does happen in spite of this wave length changes.
Cohenrence is a dynamic effect (?) and what is important in coherence
is the pulsation part that varies with time; this part is linked to
the energy of the electron and this part is not consider in the present
WP-MS static approach.

\subsection{F-MS versus AS-MS }

\label{sec:F-MSversusAS-MS } - x1 is the key parameter: AS-MS is
important only for simulations where x1 approaches $\lambda$, which
we have designated as $\lambda$-simulations - develop x1, g1, angles-wave
link, evanescent waves ?? - discussion of the singular point at k0
- but in MS the difference can more important than in free space propagator
(demander à AI pourquoi ? parce qu’il y a en plus de la doiffraction
à grand angle par les atomes !!! tout simplement)

- up to now, only few papers have performed $\lambda$-simulations
because it was not needed in TEM to reproduce the experimental images.
As resolution in TEM is approaching $\lambda$, do we need to need
to have more accuration simulation models, especially if someone is
not only interested in reproducing visually the experience but want
to use this model to reconstruct a 3D-model of the sample. - - in
experimental simulations, the wave-propagation model is not the only
limiting factors. Energy spread of the incident beam, atom vibrations,
sample and electronic drifts, lense and detector imperfections will
have their own contribution to the what is called the stobb-factors
in EM. - the purpose of these papers was not analysed all these effects
but give (i) a short picture of the different efficient numerical
approach that can be use to model the electron propagation in matter
(ii) focus on the 3 most simple and efficient ones (iii) introduce
the WP- , well known in optics but was never consider in EM, it could
be an efficient way to introduce back-propagation (that was no)

- the aim of this paper was to introduce and use efficient models
that could be useful in the on-going present development, especially
ptychography and phase reconstruction. Much more work will be neede
to see the impact the new models in phase reconstruction. (7h41)

The results point to the Fresnel free-space propagator as the main
limitation of F-MS in the high-angle and thick-specimen regimes tested
here. In the Au~{[}100{]} ODE benchmark, the discrepancy appears
first in the 135~mrad beam and only later in the central beam, after
the incorrectly propagated diffracted components have had enough thickness
to interfere back into it. In the Si~{[}111{]} CBED calculation,
the F-MS peak displacements follow the parabolic-versus-exact dispersion
error, while AS-MS remains close to WP-MS over the measured angular
range. In the thick Au~{[}100{]} CBED calculation, the F-MS amplitude
error becomes roughly an order of magnitude larger than the AS-MS
error beyond about 10~nm. These observations indicate that the principal
F-MS error is the accumulated longitudinal phase error from the parabolic
Fresnel propagator.

AS-MS removes most of the error in the F-MS method by changing only
the propagator in the propagation step. The Fresnel kernel is replaced
by the exact one-way angular-spectrum kernel, while the transmission
step and FFT structure are unchanged. WP-MS makes a further correction
by allowing the propagation kernel to depend on the local electrostatic
refractive index, so the wavelength can change with the projected
potential inside each slice. This makes WP-MS the more complete one-way
split-step model, and it is closest to the ODE reference in the Au~{[}100{]}
benchmark. Under the high-energy conditions considered here, however,
this extra local-wavelength correction is smaller than the correction
obtained by removing the Fresnel approximation. For these specimens
and energies, the main source of error is the paraxial free-space
dispersion relation rather than the absence of a potential-dependent
wavelength in the propagation step.

The silicon calculation shows that the same propagator error is not
limited to the most severe high-$Z$ case. Silicon scatters more weakly
than gold, so the CBED amplitude difference is smaller than in the
thick Au~{[}100{]} calculation. However, the F-MS peak shifts still
follow the angular dispersion error, while AS-MS remains close to
WP-MS across the measured range. A conventional image simulation can
hide part of this error if the transfer function or objective aperture
suppresses the affected Fourier components, whereas CBED and ptychographic
forward models will use the diffraction-plane intensities directly
and the error between methods will be more prominent.

In particular in a ptychographic reconstruction, the forward model
predicts detector intensities for each probe position. If the propagator
gives the wrong high-angle amplitudes, peak positions, or accumulated
phase through a thick specimen, the optimization can compensate through
the recovered potential. The thick Au~{[}100{]} CBED calculation
shows that this modeling error is much larger for F-MS than for AS-MS
after approximately 10~nm. For forward models that already use F-MS,
most of the wide-angle correction can be recovered by replacing $\exp(-i\pi k^{2}_{\perp}\Delta z/k_{0})$
with $\exp(2\pi i\Delta z[\sqrt{k^{2}_{0}-k^{2}_{\perp}}-k_{0}])$,
while keeping the same FFT structure, transmission step, and per-slice
computational complexity. The sampling estimate above shows that the
corresponding sampling penalty is only a few percent under the TEM
conditions considered here.

The comparison here is limited to high-energy, forward-scattering
calculations where the free-space dispersion error is larger than
the local wavelength correction. WP-MS should be tested separately
at lower accelerating voltages, for thicker high-$Z$ specimens, and
for slice thicknesses where the potential-dependent wavelength change
becomes a comparable source of phase error. Its main advantage in
those tests is that it provides a differentiable one-way method with
a propagation kernel that depends on the local electrostatic refractive
index, that can make use of the computational efficiency of the Fourier
transform

\section{Conclusion}

We adapted the optimised Wave Propagation Method developed in light
optics to electron microscopy, giving a wave propagation multislice
(WP-MS) method that tracks high-angle propagation and local wavelength
changes within a differentiable split-step framework, that can make
efficient use of the Fourier transform. We then compared F-MS, AS-MS,
and WP-MS in three complementary settings: an Au~{[}100{]} ODE benchmark,
a lower-$Z$ Si~{[}111{]} CBED calculation, and a thick Au~{[}100{]}
CBED calculation.

Across these tests, replacing the Fresnel kernel with the exact angular-spectrum
kernel accounts for most of the improvement over conventional multislice.
WP-MS is consistently the most complete one-way model and is closest
to the ODE reference in the Au~{[}100{]} benchmark, but its additional
accuracy over AS-MS is small under the high-energy conditions tested
here. The main practical recommendation is that for high-angle elastic
forward simulations where F-MS is still being used, AS-MS should be
the default standard, because it removes the dominant paraxial dispersion
error without changing the basic FFT-based multislice workflow.

Future work could extend the WP-MS toward a bidirectional formulation
that can include backscattering, and then benchmark it against FC-MS
in lower-energy simulations, and measure how F-MS, AS-MS, and WP-MS
affect ptychographic reconstructions across specimen thickness, accelerating
voltage, and atomic number.

\bibliographystyle{elsarticle-num}\bibliography{../sample}

\end{document}
